\chapter*{Résumé}
\addcontentsline{toc}{chapter}{Résumé}
\markboth{Résumé}{}

L'accès du citoyen marocain à l'information juridique reste difficile : les textes sont
publics, mais rédigés dans une langue dense, structurés en centaines d'articles se renvoyant
les uns aux autres, et parfois disponibles uniquement en arabe. Les assistants conversationnels
généralistes ne comblent pas ce manque : ils connaissent mal le droit marocain, peuvent inventer
des références inexistantes --- un numéro d'article plausible mais faux est particulièrement
trompeur en matière juridique --- et font transiter les questions des citoyens hors du
territoire national.

Ce stage a consisté à concevoir, réaliser et évaluer un assistant de recherche juridique
\textbf{souverain} fondé sur l'architecture RAG (\anglais{Retrieval-Augmented Generation}),
appliqué au Code du travail marocain (Loi n\textsuperscript{o}~65-99). Le système répond à une
question posée en langage naturel en s'appuyant \emph{uniquement} sur les articles réellement
retrouvés dans le corpus, cite systématiquement leurs numéros, et s'abstient explicitement
lorsque l'information demandée n'y figure pas. L'ensemble de la chaîne --- corpus, embeddings,
recherche et génération --- s'exécute en local, sans aucun appel à une interface de
programmation externe.

Le corpus a été construit à partir du PDF officiel : \NbArticles{} articles extraits, nettoyés
et enrichis de leurs métadonnées de hiérarchie, avec une régénération vérifiée identique octet
pour octet depuis la source. La recherche combine une méthode sémantique (embeddings BGE-M3,
base vectorielle Chroma) et une méthode lexicale (BM25), fusionnées par \anglais{Reciprocal Rank
Fusion}. La génération s'appuie sur un modèle Mistral~7B exécuté localement via Ollama, encadré
par deux garde-fous indépendants du modèle : un seuil d'abstention appliqué \emph{avant} l'appel
au modèle, et un contrôle vérifiant que chaque article cité figurait bien dans le contexte
fourni.

Le système a été évalué sur un jeu de référence de \NbQuestions{} questions. En configuration
hybride, la recherche atteint un Recall@5 de \RecallHybride{} et un MRR de \MrrHybride{}. Côté
réponses, \TauxCitation~\% citent au moins un article, \TauxVerif~\% des citations sont
vérifiées comme réellement fournies au modèle, et l'abstention est correcte sur
\TauxAbstention~\% des questions hors périmètre --- aucune réponse n'a été produite sur un sujet
absent du corpus. Le temps de réponse médian est de \LatTotale~secondes, réduit à
\LatAbstention~seconde lorsque le garde-fou se déclenche.

Un audit de couverture, mené une fois le système stabilisé, a montré que le jeu de référence
initial n'exerçait pas plusieurs comportements du code. Son élargissement de 37 à \NbQuestions{}
questions a \textbf{inversé} deux conclusions pourtant mesurées : la comparaison entre recherche
dense et recherche hybride, et la séparation apparente des bandes de score d'abstention. Ce
résultat méthodologique --- un banc d'évaluation est un artefact conçu, qu'il faut valider comme
le système qu'il juge --- occupe une place centrale dans ce rapport, au même titre que les
limites de mesure assumées, dont l'impossibilité de vérifier automatiquement qu'un article cité
\emph{dit réellement} ce que la réponse lui attribue.

\vspace{0.6cm}
{\color{rule}\rule{\textwidth}{0.8pt}}

\vspace{2pt}
\noindent{\sffamily\bfseries\color{navy}Mots clés~:}
RAG, souveraineté numérique, traitement automatique du langage naturel, recherche sémantique,
grands modèles de langage, droit du travail marocain, ancrage documentaire, abstention.

\vspace{2pt}
{\color{rule}\rule{\textwidth}{0.8pt}}
